\documentclass[10pt,a4paper]{article}
\usepackage{admos-article}

% ---------- Article information: edit this block ----------
\newcommand{\ArticleType}{TECHNICAL ARTICLE}
\newcommand{\ArticleSeries}{Enter Article Series Name}
\newcommand{\ArticleID}{Enter Article ID}
\newcommand{\ArticleDate}{DD.MM.YYYY}
\newcommand{\ArticleTitle}{Enter the Title of Your Article}
\newcommand{\ArticleSubtitle}{Enter a Short Subtitle or Guiding Question}

\begin{document}
	
	\pagestyle{empty}
	
	% ============================================================
	% PAGE 1 START
	% ============================================================
	
	\ArticleMetadata
	\ArticleTitleBlock
	
	Introduce the topic of the article here. Briefly explain the technical background, the problem being addressed, and why the subject is relevant. The opening paragraph should provide enough context for the reader to understand the purpose and scope of the article before moving into the detailed technical discussion.
	
	\FigurePlaceholder{
		Replace this text with a concise description of the figure. Explain what is shown and what the reader should observe or understand from it.
	}
	
	\section{Enter First Main Section Title}
	
	Use this section to introduce the first major technical topic of the article. Explain the relevant concepts, technologies, measurement principles, modeling approaches, simulation methods, or experimental procedures in sufficient detail for the intended audience.
	
	Add further paragraphs where necessary to develop the topic. When referring to figures, tables, equations, measurements, simulations, or external references, introduce them clearly in the surrounding text and explain why they are relevant.
	
	% Example for a real image:
	%
	% \begin{figure}[htbp]
		%   \centering
		%   \includegraphics[width=0.72\linewidth]{your-figure.png}
		%   \caption{
			%     Replace this text with a concise description of the figure.
			%   }
		%   \label{fig:example}
		% \end{figure}
	
	\section{Enter Second Main Section Title}
	
	Use this section for the next major aspect of the article. Depending on the subject, this section could present measurement results, simulation results, model development, parameter extraction, device behavior, experimental observations, comparison of different methods, or another relevant technical topic.
	
	Explain the results or observations rather than simply presenting them. Where appropriate, discuss trends, differences, important operating conditions, limitations, or unexpected behavior.
	
	The example table below can be replaced with article-specific data. Adjust the column headings, entries, and number of rows according to the information that needs to be presented.
	
	\begin{table}[htbp]
		
		\centering
		
		\begin{tabular}{@{}ll@{}}
			
			\toprule
			
			\textbf{Enter Column Heading}
			&
			\textbf{Enter Column Heading}
			\\
			
			\midrule
			
			Enter value & Enter value \\
			Enter value & Enter value \\
			Enter value & Enter value \\
			Enter value & Enter value \\
			
			\bottomrule
			
		\end{tabular}
		
		\caption{
			Replace this text with a concise description of the information presented in the table.
		}
		
	\end{table}
	
	\section{Enter Additional Section Title}
	
	Use additional sections when required to organize the technical discussion. Each section should focus on one clearly defined topic and should follow logically from the preceding section.
	
	Where appropriate, discuss important observations, assumptions, limitations, comparisons, or practical implications. Figures and tables should support the technical discussion and should be explained in the surrounding text.
	
	\subsection{Enter Subsection Title if Required}
	
	Use subsections when a main section contains several closely related topics that benefit from additional organization.
	
	Keep subsection titles concise and descriptive. The text should remain focused on the topic introduced by the main section.
	
	\section{Conclusion}
	
	Summarize the main technical findings and key messages of the article here. Focus on the most important conclusions rather than repeating the complete discussion.
	
	Where appropriate, mention the practical relevance of the results, important limitations, remaining challenges, possible improvements, or topics that could be investigated in future work.
	
	% ============================================================
	% REFERENCES
	% ============================================================
	
	\begin{thebibliography}{9}
		
		\bibitem{reference1}
		A. Author,
		``Enter the title of the referenced journal or conference article,''
		\textit{Journal or Conference Name},
		vol. X,
		no. X,
		pp. XX--XX,
		Year.
		
		\bibitem{reference2}
		B. Author,
		\textit{Enter the Title of the Referenced Book}.
		City, Country:
		Publisher,
		Year.
		
		\bibitem{reference3}
		Organization or Author,
		``Enter the title of the referenced standard, website, technical report, or documentation,''
		Year.
		
	\end{thebibliography}
	
\end{document}